\makeatletter
\def\input@path{{../styles/}{../../styles/}{../../../styles/}{../}{../../}{../../../}}
\makeatother

\documentclass{ee102_pset}
% macros.tex - Course meta information
\renewcommand{\course}{EE 102} % with a space
\renewcommand{\coursetitle}{Signal Processing and Linear Systems}
\renewcommand{\instructor}{Ayush Pandey}
\renewcommand{\student}{Name: }

\renewcommand{\semester}{Fall 2025}
\date{\semester} % this sets the LaTeX date field safely

% Problem set number
\renewcommand{\psetnum}{1}

% Release / due — use renewcommand because package provides empties
\renewcommand{\releasedate}{August 20, 2025}
\renewcommand{\duedate}{August 27, 2025}

\input{preamble.tex}
\usepackage{listings}

\author{\rule{3cm}{0.4pt}}
\submitdate{\rule{3cm}{0.4pt}}
\problemset{Problem Set \#11: Filter Design}
\renewcommand{\psetnum}{11}
\renewcommand{\releasedate}{November 11, 2026}
\renewcommand{\duedate}{Nov 18, 2026 at 11:59 PM}
\shorttitle{Problem Set \#11}
\hypersetup{pdftitle={EE 102 Problem Set 11: Filter Design}}

\begin{document}

\problem{1} An underdamped series RLC circuit is excited by an input at time \(t_0\). The circuit outputs a voltage that is a damped oscillation at its natural frequency. This output signal is modeled as
\[
x(t)=A\,e^{-\alpha\,(t-t_0)}\cos\!\big(\omega_0\,(t-t_0)\big)\,u(t-t_0),
\qquad
\alpha>0,\ \omega_0>0,
\]
where \(u(\cdot)\) is the unit step. 

\problempart \textbf{(3 points)}
Write two time-domain properties of \(x(t)\) that you can observe directly from the signal equation or by graphing it. You may graph the signal using any software (Python, MATLAB, Desmos) or by hand.

\problempart \textbf{(2 points)} Write a signal $\bar{x}(t)$ that when shifted by \(t_0\) gives \(x(t)\).

\problempart \textbf{(10 points)}
Compute the Fourier transform \(X(\omega)\) of \(x(t)\).

Hint: Use Euler's formula to express the cosine as a sum of complex exponentials --- don't solve integration by parts!

\problempart \textbf{(5 points)} Compute the Fourier transform $\bar X(\omega)$ of $\bar x(t)$. Using the time-shift property, prove that \(X(\omega) = e^{-j\omega t_0} \bar{X}(\omega)\) to verify your result.


% \[ X(\omega)=\frac{A}{2}\,e^{-j\omega t_0}\!\left(\frac{1}{\alpha+j(\omega-\omega_0)}+\frac{1}{\alpha+j(\omega+\omega_0)}\right).
% \]

% \problempart \textbf{(5 points)} Sketch \(|X(\omega)|\). Indicate the spectral peaks (the frequency peaks) at \(\pm\omega_0\).

\problempart \textbf{(5 points)}
Define an ideal low-pass filter as follows:
\[
H_{\mathrm{LP}}(\omega)=
\begin{cases}
1,& |\omega|\le \Omega_c,\\
0,& |\omega|>\Omega_c.
\end{cases}
\]
Sketch \(H_{\mathrm{LP}}(\omega)\) for positive and negative \(\omega\). Label \(\pm\Omega_c\) on your graph.

\problempart \textbf{(5 points)}
By using the inverse Fourier transform, find the impulse response of the low-pass filter. 

% \[
% h_{\mathrm{LP}}(t)=\frac{1}{2\pi}\!\int_{-\Omega_c}^{\Omega_c} e^{j\omega t}\,d\omega
% =\frac{\sin(\Omega_c t)}{\pi t}.
% \]

\textbf{Hint:} use even symmetry of \(H_{\mathrm{LP}}(\omega)\). Solved example 4.5 in Oppenheim \& Willsky 2nd Edition has the answer. It is OK to just use the result and approach from the book.

\problempart \textbf{(10 points)}
Now that we have the impulse response of the system, you know that you can find the output of the system, \(y(t)\), for any input \(x(t)\) by using the convolution integral. But the convolution theorem we discussed in class gave us a much easier way to do this. That is, we can find the frequency domain representation of the output by simple multiplication:
\[
Y(\omega)=X(\omega)\,H_{\mathrm{LP}}(\omega).
\]
Compute \(Y(\omega)\) and discuss its properties.

\problempart \textbf{(10 points)}
But still, to prove the point that frequency domain actually made the computation easier, let us try to attempt the same problem in time-domain. Write the convolution integral
\[
y(t)=(x*h_{\mathrm{LP}})(t)=\int_{-\infty}^{\infty} x(\tau)\,h_{\mathrm{LP}}(t-\tau)\,d\tau.
\]
Solve the convolution integral to compute $y(t)$ in time-domain.
Compute the inverse Fourier transform of \(Y(\omega)\) to obtain \(y(t)\) and verify your result. 

\textbf{\color{blue}Important:} You will not be able to solve the convolution integral in closed-form. Instead, express your answer as an integral expression that can be evaluated numerically. Things can be quite difficult in time-domain!

\problempart \textbf{(15 points)}
Using Parseval's theorem, we can define a new type of energy in the frequency domain that is band-limited to \(|\omega|\le \Omega_c\):
\[
E_{\mathrm{LP}}=\frac{1}{2\pi}\int_{-\Omega_c}^{\Omega_c} |X(\omega)|^2\,d\omega,
\]
this is the energy contained in the signal \(x(t)\) that lies within the passband of the low-pass filter.

Compute \(E_{\mathrm{LP}}\) in terms of \(A,\alpha,\omega_0,\Omega_c\).
Write the exact integral using your expression for \(X(\omega)\), then evaluate it numerically for the parameter values you choose in Problem 2. You do not need a closed-form answer.

\noindent\begin{minipage}{\linewidth}
\textbf{Python hint:} Define \texttt{X(w)} using your Fourier transform from part (c), with numerical values for the parameters and \texttt{Omega\_c}.
\begin{lstlisting}[language=Python,basicstyle=\ttfamily\small,frame=single,breaklines=true,columns=fullflexible]
import numpy as np
from scipy.integrate import quad

E_LP, error = quad(lambda w: abs(X(w))**2,
                   -Omega_c, Omega_c)
E_LP = E_LP / (2 * np.pi)
print(E_LP)
\end{lstlisting}
\end{minipage}

\noindent\begin{minipage}{\linewidth}
\textbf{MATLAB hint:} Define a vectorized function handle \texttt{X} using your Fourier transform from part (c), with numerical values for the parameters and \texttt{Omega\_c}. Use elementwise operations in \texttt{X}.
\begin{lstlisting}[language=Matlab,basicstyle=\ttfamily\small,frame=single,breaklines=true,columns=fullflexible]
E_LP = integral(@(w) abs(X(w)).^2, ...
                -Omega_c, Omega_c) / (2*pi);
disp(E_LP)
\end{lstlisting}
\end{minipage}

\problempart \textbf{(10 points)}
Explore other possible filters in frequency domain: low-pass, band-pass, and high-pass. Comment on which one might better isolate the ring-down mode in the RLC circuit. Refer to sketches, \(h(t)\), and band-energy calculations. You don't need to solve it out completely but make convincing arguments using Fourier analysis properties.


\problem{2} Validate Problem 1 using computer programs. It is \textbf{highly recommended that you start early with the homework so that you can get help during your discussion section}.

With a sampling rate \(F_s\) Hz, we have a time interval of \(\Delta t=1/F_s\) and we can write the continuous-time signal above by substituting \(t=n\Delta t\) as a discrete-time signal:
\[
x[n]=A\,e^{-\alpha(n-n_0)\Delta t}\cos\!\big(\omega_0 (n-n_0)\Delta t\big)\,u[n-n_0].
\]
Here \(n_0=t_0/\Delta t\).

\problempart \textbf{(5 points)}
Choose parameters \((A,\alpha,\omega_0,t_0)\) and \(F_s\). Implement the equation above in a function. Then, plot \(x[n]\) versus \(n\Delta t\). Label your graph.
Make sure to choose $F_s$ to be greater than \(2 \cdot \frac{\omega_0}{2\pi}\). You will learn why this is important next week. 

For a simple start, set $F_s = 2000$ Hz and ensure that the condition above is satisfied.

\problempart \textbf{(10 points)}
Re-use the previous homework assignment code for DFT implementation to compute the DFT with length \(N\) for the signal above (it is also OK to use FFT libraries). 

Helpful starting point:

Set $N = 2048$ and $L = 8192$. If your generated \(x[n]\) has length \(<N\), zero-pad it to length \(N\). If it is longer, keep the first \(N\) samples. Then zero-pad from \(N\) to \(L\) before taking the DFT. Remember that with DFT, the frequency is sampled at discrete frequencies. Use the following frequency grid:
\[
\omega_k=\frac{2\pi F_s}{L}
\begin{cases}
k,&0\le k<L/2,\\
k-L,&L/2\le k<L.
\end{cases}
\]
The index corresponding to \(\omega_0\) will be 
\[
k_0=\operatorname{round}\!\left(\frac{\omega_0\,L}{2\pi F_s}\right).
\]

Plot \(|X[k]|\) versus discrete frequency. Indicate indices corresponding to \(\pm\omega_0\). You may need to center the spectrum for better visualization (using \texttt{fftshift}).

\problempart \textbf{(10 points)}
Implement the ideal low-pass filter $H[k]$. Compute \(y[n]\) using inverse DFT. Plot the filtered signal \(y[n]\).

The discrete-time DFT mask is
\[
H[k]=\begin{cases}
1,& |\omega_k|\le \Omega_c,\\
0,& \text{otherwise}.
\end{cases}
\]
You may start by setting \(\Omega_c = 0.8\,\omega_0\).

Compute \(Y[k]=X[k]\cdot H[k]\). Compute \(y[n]\) by the \(L\)-point inverse DFT. Plot \(\operatorname{Re}\{y[n]\}\) versus \(n\Delta t\).

\textbf{\color{blue}Optional exploration (not graded):}

\problempart Implement at least one other filter type (high-pass or band-pass, or something else of your choice). Provide the DFT mask, impulse response, and plot the filtered signal. Comment on how well the filter isolates the ring-down mode in the RLC circuit.


\problempart Estimate band-limited input energy by summing \(|X[k]|^2\) over passband indices to verify your answer in Problem 1. Find out how well your filter isolates the ringing mode for your chosen parameters. Experiment with different parameters.

\end{document}
